\documentclass[10pt,a4paper]{article}

\usepackage[T1]{fontenc}
\usepackage[utf8]{inputenc}
\usepackage[margin=12mm,top=27mm,bottom=13mm]{geometry}
\usepackage[table]{xcolor}
\usepackage{array}
\usepackage{booktabs}
\usepackage{tabularx}
\usepackage{multicol}
\usepackage{fancyhdr}
\usepackage{microtype}
\usepackage{tikz}
\usepackage{hyperref}

\definecolor{Navy}{HTML}{15324A}
\definecolor{Teal}{HTML}{168C8C}
\definecolor{PaleTeal}{HTML}{E8F4F3}
\definecolor{Ink}{HTML}{23313B}
\definecolor{Muted}{HTML}{60717C}
\definecolor{Grid}{HTML}{C8D3D8}
\definecolor{Pale}{HTML}{F4F7F8}
\definecolor{BlueGray}{HTML}{8BA6B1}

\hypersetup{
  colorlinks=true,
  linkcolor=Teal,
  urlcolor=Teal,
  pdftitle={Problem Definition Report - EVS 2017},
  pdfauthor={Student report},
  pdfsubject={Data understanding for a semester project}
}
\urlstyle{same}

\renewcommand{\familydefault}{\sfdefault}
\setlength{\parindent}{0pt}
\setlength{\parskip}{5.5pt plus 0.6pt minus 0.4pt}
\setlength{\columnsep}{5mm}
\setlength{\columnseprule}{0pt}
\raggedcolumns
\color{Ink}

\newcommand{\StudentName}{\rule{35mm}{0.3pt}}
\newcommand{\DatasetFooter}{Joint EVS/WVS 2017--2022 v5.0.0 - EVS analytical subset}
\newcommand{\ReportHeader}{%
  \colorbox{Navy}{%
    \parbox{\dimexpr\textwidth-2\fboxsep\relax}{%
      \color{white}\vspace{1.3mm}%
      {\fontsize{14}{15}\selectfont\bfseries Problem Definition Report - Part I: Data}\\[0.6mm]
      {\fontsize{8.4}{9.2}\selectfont EVS 2017: explaining life satisfaction across countries}\\[1mm]
      {\fontsize{6.3}{7}\selectfont\color{white!82}Student: \StudentName\quad |\quad 4EK481 / 4EK666\quad |\quad 4 October 2026}%
      \vspace{1.2mm}%
    }%
  }%
}

\pagestyle{fancy}
\fancyhf{}
\setlength{\headheight}{20mm}
\fancyhead[C]{\ReportHeader}
\renewcommand{\headrulewidth}{1.6pt}
\renewcommand{\headrule}{\hbox to\headwidth{\color{Teal}\leaders\hrule height \headrulewidth\hfill}}
\fancyfoot[L]{\fontsize{6.2}{7}\selectfont\color{Muted}\DatasetFooter}
\fancyfoot[R]{\fontsize{6.2}{7}\selectfont\color{Muted}Page \thepage}
\renewcommand{\footrulewidth}{0.35pt}

\newcommand{\ReportSection}[2]{%
  \vspace{3pt}%
  {\fontsize{9.5}{10.5}\selectfont\bfseries\color{Navy}#1. #2}\par
  \vspace{1pt}%
}

\newcommand{\EvidenceHeading}{%
  \vspace{3pt}%
  {\fontsize{9.5}{10.5}\selectfont\bfseries\color{Navy}Methods, evidence and references}\par
  \vspace{1pt}%
}

\newcommand{\KeyBox}[1]{%
  \begingroup
  \setlength{\fboxsep}{5pt}%
  \noindent\fcolorbox{Teal}{PaleTeal}{\parbox{\dimexpr\linewidth-2\fboxsep-2\fboxrule\relax}{#1}}%
  \endgroup
}

\newcolumntype{Y}{>{\raggedright\arraybackslash}X}
\newcolumntype{L}[1]{>{\raggedright\arraybackslash}p{#1}}
\newcommand{\TableHead}[1]{\textcolor{white}{\bfseries #1}}

\newenvironment{ReportTable}[1]
  {\rowcolors{2}{Pale}{white}\fontsize{6.35}{7.45}\selectfont%
   \setlength{\tabcolsep}{2.8pt}\renewcommand{\arraystretch}{1.40}%
   \tabularx{\linewidth}{#1}\rowcolor{Navy}\color{white}}
  {\endtabularx}

\AtEndDocument{%
  \ifnum\value{page}=2\relax
  \else
    \PackageError{evs-report}{Expected exactly two pages, got \thepage}{}
  \fi
}

\begin{document}
\fontsize{7.75}{9.25}\selectfont

\begin{multicols}{2}

\KeyBox{%
  \textbf{Research question}\par
  Which demographic, socioeconomic and social characteristics are associated with life satisfaction in EVS 2017, and how stable are those relationships across countries?\par
  \vspace{2.5pt}\hrule height 0.3pt\vspace{2.5pt}
  \textbf{Analytical scope}\par
  59,438 respondents $\mid$ 36 countries $\mid$ 226 available variables $\mid$ 2017--2021
}

\ReportSection{1}{Provenance and purpose}
\textbf{Dataset.} Joint EVS/WVS 2017--2022 Dataset, ZA7505, version 5.0.0; this report uses only European Values Study records (\texttt{study = 1}) from the joint common-core file. The original producers are the European Values Study and World Values Survey; GESIS archives and publishes the joint file. Persistent source: \href{https://doi.org/10.4232/1.14320}{doi:10.4232/1.14320}.[1]

\textbf{Collection and domain.} EVS planned country surveys under common methodological guidance to study values, beliefs, attitudes and well-being. National probability samples interviewed adults in private households during 2017--2021; EVS and WVS then harmonised common questionnaire items.[2] It is interesting for this project because it connects individual circumstances and social context with life satisfaction across diverse European societies.

\textbf{Access/terms.} Data and documentation are available free of charge after GESIS registration, are governed by GESIS Data Usage Terms, and should be cited.[3] No open-content licence label is asserted in the supplied variable report.

\ReportSection{2}{Size and computational profile}
\begin{ReportTable}{L{0.44\linewidth}Y}
\TableHead{Quantity} & \TableHead{Result}\\
Official joint dataset & 156,658 respondents from EVS and WVS\\
Selected EVS subset ($n$) & 59,438 respondents; 37.94\% of the official joint dataset\\
Features $p$ & 226 non-empty columns (231 raw; 5 WVS-only columns removed)\\
Files / format & 1 official CSV / 1 table\\
Raw joint CSV size & 104.3 MiB\\
EVS working file & 3.8 MiB Parquet (40.4 MiB CSV-equivalent)\\
Dense 8-byte baseline & 104.8 MiB plus string/software overhead\\
\end{ReportTable}

The official joint file is substantially larger: it contains 156,658 EVS and WVS respondents, while the selected 59,438 EVS records represent 37.94\% of all rows. One observation is one respondent; one feature is one delivered column before excluding IDs, protocol fields, weights and $Y$. The data fit comfortably in laptop memory. Size is mainly in $n$, but the greater difficulty is hierarchical dependence and mixed measurement scales. With $n/p \approx 263$ for the 226-column subset, $p$ is small relative to $n$; regularisation may still matter after categorical expansion and with correlated survey items.

\ReportSection{3}{Observational structure: what is one row?}
The data are \textbf{hierarchical pooled cross-sections}, not a person-level panel: respondents are nested within 36 countries, survey years (2017--2021), fieldwork designs and interview modes. All 59,438 \texttt{uniqid} values are unique. Rows are distinct people but not fully independent: respondents within a country and fieldwork context share institutions, shocks, translations and survey procedures.

A random respondent split would put the same countries and contexts in train and test sets and can overstate transfer performance. Use grouped country holdouts to assess new-country generalisation; where possible, use later fieldwork as a temporal test. Avoid identifiers, fieldwork dates and protocol variables as ordinary predictors because they may leak context.

\textbf{Coverage.} AL, AM, AT, AZ, BA, BG, BY, CH, CZ, DE, DK, EE, ES, FI, FR, GB, GE, HR, HU, IS, IT, LT, LV, ME, MK, NL, NO, PL, PT, RO, RS, RU, SE, SI, SK, UA. Country samples range from 1,003 to 3,362.

\columnbreak

\ReportSection{4}{Variables, support and representation}
The table is a representative shortlist of important variables, not the complete set of 226 available columns.

\begin{ReportTable}{L{0.32\linewidth}L{0.43\linewidth}Y}
\TableHead{Role / variable} & \TableHead{Type and support} & \TableHead{Missing}\\
$Y$: A170 life satisfaction & Ordinal 1--10; numeric treatment is an assumption & 316 (0.53\%)\\
A008 happiness & Ordinal categorical & 731 (1.23\%)\\
A009 health & Ordinal categorical & 134 (0.23\%)\\
A165 trust & Binary categorical & 1,355 (2.28\%)\\
A173 control & Ordinal 1--10 & 888 (1.49\%)\\
X001 / X003 sex / age & Nominal / integer; age 82 means 82+ & 26 / 325\\
X025R / X028 education / work & Ordinal / nominal & 472 / 411\\
X047E\_EVS5 income & Ordinal household-income decile & 8,740 (14.70\%)\\
\end{ReportTable}

The remaining columns cover administrative IDs, survey protocol and weights; perceptions of life; work; family and household; national identity; politics and society; religion and morality; security; and respondent, parent and partner demographics. Many administrative fields are not suitable predictors. The model receives integer-coded survey responses, not raw speech or behaviour; wording, response nuance and within-category variation are already discarded. Convert negative missing codes before fitting, encode nominal variables, and preserve ordinal order where the model supports it. Standardise only scale-sensitive numeric inputs; logs are not generally appropriate for bounded categories. Euclidean distance on raw category codes is not meaningful, and many attitude items may be redundant or noisy.

\vspace{2pt}
{\fontsize{6.7}{7.8}\selectfont\bfseries\color{Navy}A170: life satisfaction (valid $n = 59{,}122$)}\hfill
{\fontsize{6.2}{7}\selectfont\color{Teal}72.5\% chose 7--10}\par
\centering
\begin{tikzpicture}[x=0.67cm,y=0.00014cm]
  \foreach \x/\v in {0/1089,1/731,2/1557,3/2047,4/5370,5/5458}{
    \fill[BlueGray] (\x,0) rectangle +(0.50,\v);
  }
  \foreach \x/\v in {6/9913,7/14594,8/9126,9/9237}{
    \fill[Teal] (\x,0) rectangle +(0.50,\v);
  }
  \foreach \x/\lab in {0/1,1/2,2/3,3/4,4/5,5/6,6/7,7/8,8/9,9/10}{
    \node[font=\fontsize{5.5}{6}\selectfont,text=Muted] at (\x+0.25,-1650) {\lab};
  }
\end{tikzpicture}
\raggedright

{\fontsize{6.15}{7.2}\selectfont\color{Muted}\textit{Figure 1.} A170 is concentrated at high values (mean 7.37; median 8), so performance should be compared with a simple central-value baseline.}

\end{multicols}

\newpage
\begin{multicols}{2}

\ReportSection{5}{Missingness and incomplete observation}
There are no blank CSV cells; incompleteness is encoded explicitly. Across all selected cells: \textbf{-1 do not know} 303,010 (2.21\%); \textbf{-2 no answer} 112,462 (0.82\%); \textbf{-3 not applicable} 570,604 (4.16\%); \textbf{-4 not asked} 412,640 (3.01\%); \textbf{-5 other missing} 5,956 (0.04\%). These categories must remain distinguishable during preparation.

\begin{ReportTable}{L{0.34\linewidth}L{0.22\linewidth}Y}
\TableHead{Key variable / meaning} & \TableHead{Valid} & \TableHead{Missing and interpretation}\\
\shortstack[l]{\texttt{A170}\\Life satisfaction} & 59,122 (99.47\%) & \shortstack[l]{199 (-1: do not know)\\117 (-2: no answer)}\\
\shortstack[l]{\texttt{X047E\_EVS5}\\Net household-income decile} & 50,698 (85.30\%) & \shortstack[l]{2,513 (-1: do not know)\\5,011 (-2: no answer)\\1,216 (-5: other missing)}\\
\shortstack[l]{\texttt{A165}\\Most people can be trusted} & 58,083 (97.72\%) & \shortstack[l]{1,012 (-1: do not know)\\343 (-2: no answer)}\\
\end{ReportTable}

Missingness is unlikely to be globally MCAR. \textbf{-4} is structural by questionnaire/country; MAR is plausible conditional on country, mode, year and demographics; refusals or uncertainty on income may be MNAR. Do not impute at this stage. Age, children and household size are upper-censored rather than missing: 82 means 82+, 5 children means 5+, and household size 6 means 6+. People outside private households are absent from the sample (unit truncation/undercoverage).

\ReportSection{6}{Outliers, measurement and data quality}
\textbf{Checks.} No duplicate IDs and no values outside documented support were found for the target or selected predictors. A concrete unusual case is the top-coded age category: 1,666 respondents (2.80\%) are recorded as 82, a genuine censored group rather than necessarily an error. Likewise, 1,577 respondents report the 5+ children category and 3,125 the 6+ household-size category. Negative values are documented missing codes, not outliers.

Survey weights can look extreme but are purposeful: \texttt{gwght} ranges 0.006--13.552, \texttt{pwght} 0.015--6.639, and \texttt{wght\_eq1000} 0.297--0.997. Inspect influence and use the weight matching the estimand rather than deleting large weights mechanically. Country labels and units are harmonised, but source-specific fields such as \texttt{X047\_WVS7} are structurally unavailable for EVS and must be excluded.

Life satisfaction, happiness, health, trust, perceived control and income are self-reported and vulnerable to recall, social-desirability, translation, mode and cultural response-style error. Important unmeasured factors include recent life events, detailed mental health, wealth, personality, local prices and country-time shocks.

\ReportSection{7}{Sample, population and future use}
\textbf{Population of interest.} Adults aged 18+ living in private households in the 36 EVS countries during the wave, with cautious extension to later adults in those countries. Country surveys used single- or multistage probability samples under EVS guidance, so the design is much stronger than convenience sampling, but nonresponse, household-frame exclusions and unequal country sample sizes can bias pooled results. Survey weights adjust selected margins, not all selection or measurement bias.

Coverage is broad geographically but not all of Europe, and each country contributes roughly 1,000--3,400 people; rare demographic combinations and institutional contexts may be sparse. Extrapolation to excluded countries, people outside private households, very small subgroups or future crises is risky.

Future data may shift with country mix, ageing, migration, economic or political conditions, interview mode, translation and post-pandemic attitudes. Both covariate and concept shift are plausible. The least stable relationship is expected across countries and modes, and between pre-COVID and 2020--2021 fieldwork. Report country-specific performance and uncertainty rather than only one pooled score.

\columnbreak

\ReportSection{8}{Signal, noise and first modelling intuition}
The objective is to \textbf{predict and explain variation in A170}, while treating findings as associations rather than causal effects. Expected signal: subjective health, happiness, perceived control, employment, education, income, age, household composition and country context. Expected irreducible noise: temporary mood, omitted personal events, response-scale interpretation and local shocks.

A170 is ordinal, not a rare event: all ten scores occur, although 72.51\% of valid responses are 7--10. Compare ordinal regression and a carefully justified numeric model with a simple baseline; include country effects and grouped validation. Closely related subjective items may be strong predictors but can blur explanation and outcome definition. \textbf{Main anticipated difficulty:} cross-country and cross-mode heterogeneity combined with survey dependence, not computation.

\EvidenceHeading
All counts in this report were computed from the supplied immutable CSV after selecting \texttt{study = 1}; the five all-missing WVS-only columns were removed only in the derived EVS Parquet file. Missing codes were not imputed. Variable meanings and supports were checked against the supplied ZA7505 variable report.[4]

{\fontsize{6.35}{7.5}\selectfont\color{Muted}
[1] EVS/WVS (2024), \textit{Joint EVS/WVS 2017--2022 Dataset}, ZA7505 v5.0.0, GESIS: \href{https://doi.org/10.4232/1.14320}{doi.org/10.4232/1.14320}.\par
[2] EVS, joint-dataset description and composition: \href{https://europeanvaluesstudy.eu/surveys/evs-wvs-collaborations/joint-evs-wvs/}{europeanvaluesstudy.eu/.../joint-evs-wvs/}.\par
[3] EVS/GESIS data access and usage terms: \href{https://europeanvaluesstudy.eu/surveys/data-access/}{europeanvaluesstudy.eu/surveys/data-access/}.\par
[4] \textit{Joint EVS/WVS 2017--2022 Variable Report}, ZA7505 v5.0.0 (supplied documentation).
}

\end{multicols}
\end{document}
